\documentclass[a4paper]{article}

\usepackage{graphicx}
\usepackage{amsfonts}
\usepackage{geometry}
\usepackage[utf8]{inputenc}
\usepackage{array}
\usepackage{calc}
\usepackage{layout}

\usepackage{listings}
\geometry{
	a4paper,
	total={170mm,257mm},
	left=15mm,
	top=15mm,
}

\title{\begin{center}
		\setlength{\tabcolsep}{0pt}
		\begin{tabular}{>{\raggedleft}m{3cm}>{\centering}m{\textwidth - 5cm\relax}>{\raggedright}m{3cm}}
			\includegraphics[width=\linewidth]{poliba.png}%
			&%
			\textbf{ } \\[5pt]%
			\textbf{\ }%
			&%
			\includegraphics[width=\linewidth]{dei.png} %
		\end{tabular}
	\end{center}
	\textbf{Relazione attività di laboratorio\\}
	\textbf{\large Esercitazione 2\\- Corso di Fondamenti delle Misurazione -\\ }
	{\normalsize 
		- Ingegneria Informatica e dell'Automazione -
}}
\author{A cura di:\\
	Laura Loperfino\\
	Andrea Lops \\
	Teresa Pantone\\
	Paolo Rotolo\\
	\\
	Prof. Ing. Filippo Attivissimo
}

\date{19/12/2018}


\begin{document}

\maketitle

\lstset{language=Matlab}

\section {Incertezza su grandezze dimensionali}
%TODO:  aggiungere commento @paolo 
\subsection{Misure con centimetro estensibile}
L'unica incertezza che interressa le misure effettuate con il centimetro è l'incertezza assoluta di lettura, ovvero il massimo errore che si commette nell'apprezzare una frazione di divisione

\begin{figure}[htp]
	\centering
	\includegraphics[scale=0.80]{"fig1".jpeg}
	\caption{Centimetro estensibile}
	\label{}
\end{figure}
\begin{Large} 
	\begin{equation}
	\lambda_{min}=\frac{1}{2}*\frac{1}{10} div= \frac{1}{20}div 
	\end{equation}
\end{Large}
\begin{Large} 
	\begin{equation}
	U_{\lambda_{min}}=\frac{1}{2}*\frac{1}{20} = \frac{1}{40} 
	\end{equation}
\end{Large}
\begin{Large} 
	\begin{equation}
	\Delta_x=2U_{\lambda_{min}} = 0.5 mm 
	\end{equation}
\end{Large}
Si è calcolata lincertezza delle relative misure: 
\begin{Large} 
	\begin{equation}
		l_a = 1180\pm0.5mm
	 \end{equation}
	 \begin{equation}
		l_v = 1000\pm0.5mm
	 \end{equation}
\end{Large}
\subsection{Misure col calibro digitale}
Per quanto riguarda invece le misure effettuate col calibro digitale si ha: 
\begin{Large} 
	\begin{equation}
		U_{calibro} = 0.01 + |0.02| = 0.03mm
	 \end{equation}
\end{Large}
Si è calcolata lincertezza delle relative misure: 
\begin{Large} 
	\begin{equation}
		w = 30.08\pm0.03mm
	 \end{equation}
	 \begin{equation}
		h = 3.05\pm0.03mm
	 \end{equation}
\end{Large}
\subsection{Misure di superfici}
Successivamente sono state effettuate le misure indirette con i dati acquisiti:
\begin{Large} 
	\begin{equation}
		S = h*w = 91.744 mm^2
	 \end{equation}
\end{Large}
Con relativa incertezza:
\begin{Large} 
	\begin{equation}
		U_{S} =S(\frac{U_h}{h}+\frac{U_w}{w})
	 \end{equation}
\end{Large}
E quindi: 
\begin{Large} 
	\begin{equation}
		S =91.7\pm 1.0mm^2
	 \end{equation}
\end{Large}
Invece:
\begin{Large} 
	\begin{equation}
		S_L= 2l_a(h+w) = 78186.8 mm^2
	 \end{equation}
\end{Large}
Con relativa incertezza:
\begin{Large} 
	\begin{equation}
		U_{S_L} =\frac{dS_L}{dl_a}U_{l_a}+\frac{dS_L}{dw}U_w+\frac{dS_L}{dh}U_h=274.12 mm^2
	 \end{equation}
\end{Large}
E quindi: 
\begin{Large} 
	\begin{equation}
		S_L =78200\pm 300mm^2
	 \end{equation}
\end{Large}
\subsection{Stima del valore del provino in rame}
Usando la formula forniteci e i risultati ottenuti precedentemente è stato possibile misurare:
\begin{Large}
	\begin{equation}
		R_{X_{m}} =\frac{\varphi*l_v}{S}= 0.1918 m\Omega
	 \end{equation}
\end{Large}
Con la relativa incertezza
\begin{Large}
	\begin{equation}
		U_{R_{X_{m}}} = \varphi\frac{d\frac{\varphi*l_v}{s}}{dl_v}U{l_v}+\varphi\frac{d\frac{\varphi*l_v}{s}}{dS}U{S}
	 \end{equation}
	 \begin{equation}
		R_{X_{m}} =0.192\pm 0.004 m \Omega
	 \end{equation}
\end{Large}


\section{Dimensionamento $I_{max}$}
Con il seguente programma in Matlab è stato possibile definire $I_{max}$

\lstinputlisting{main.m}

Con risultato: 
\begin{Large}
	\begin{equation}
		I_{max} = 20.82 A
	 \end{equation}
\end{Large}


\section {Incertezza su grandezze elettriche, metodo voltamperometrico}
Per effettuate le misure seguenti sono stati adoperati: il multimetro da banco \emph{Agilent 34410} e il multimetro palmare \emph{U/U1253B}.
\subsection{Misura diretta della resistenza mediante due morsetti}
Adoperando due morsetti si è arrivati alla misurazione della seguente resistenza:
\begin{Large}
	\begin{equation}
		R_{X2W}= 0.189 \Omega
	\end{equation}
	\begin{equation}
		U_{R_{X2W}} = [\pm 0.01\%rdg \pm 0.004\% FSO]
	 \end{equation}
	 \begin{equation}
		R_{X2W} = 0.189 \pm 0.004 \Omega
	 \end{equation}
\end{Large}
Purtroppo però la misura è \emph{errata} a causa delle resistenze dei due puntali, infatti esse come tutta la cavetteria hanno una resistenza interna.
Il problema è fortunatamnete di facile risoluzione, compensando i puntali stessi. Ed ecco che si ottiene una misurazione più precisa.

\begin{Large}
	\begin{equation}
		R_{X2W_{2}} = 0.007 \Omega
	 \end{equation}
	\begin{equation}
		U_{R_{X2W_{2}}} = [\pm 0.01\%rdg \pm 0.004\% FSO]
	 \end{equation}
	 \begin{equation}
		R_{X2W_{2}} = 0.007 \pm 0.004 \Omega
	 \end{equation}
\end{Large}
\subsection{Misura della resistenza mediante quanto morsetti}
Quello che si è andato a misurare con quattro puntali è solamente la \emph{caduta di tensione} nel resistore; infatti il flusso di corrente misurato nei puntali è pressochè nullo.
\begin{Large}
	\begin{equation}
		R_{X4W}= 0.1909 \Omega
	\end{equation}
	\begin{equation}
		U_{R_{X4W}} = [\pm 0.01\%rdg \pm 0.004\% FSO]
	 \end{equation}
	 \begin{equation}
		R_{X4W} = 0.1909\pm 0.006 \Omega
	 \end{equation}
\end{Large}
\subsection{Misura di resistenza con il metodo amperometrico}
Per ottenere la resistenza $R_{X_{VA}}$ si è andati a calcolare prima $I_{VA}$ e poi $V_{VA}$
\begin{Large}
	\begin{equation}
		I_{VA}= 4.9433 A
	\end{equation}
	\begin{equation}
		U_{I_{VA}} = [\pm 0.3\%rdg \pm 10\% FSO]
	 \end{equation}
	 \begin{equation}
		I_{VA} = 4.9 \pm 0.5 A
	 \end{equation}
\end{Large}

\begin{Large}
	\begin{equation}
		V_{VA}= 0.96 mV
	\end{equation}
	\begin{equation}
		U_{V_{VA}} = [\pm 0.005\%rdg \pm 0.0035\% FSO]
	 \end{equation}
	 \begin{equation}
		V_{VA} = 0.96 \pm 0.004 mV
	 \end{equation}
\end{Large}
E successivamente la resistenza:
\begin{Large}
	\begin{equation}
		R_{X_{VA}}= \frac{V_{VA}}{I_{VA}} =0.000194 \Omega
	\end{equation}
	\begin{equation}
		U_{R_{X_{VA}}} = (\frac{U_{V_{VA}}}{V_{VA}}+\frac{U_{I_{VA}}}{I_{VA}})R_{X_{VA}}
	 \end{equation}
	 \begin{equation}
		R_X{_{VA}} = 0.00019\pm 0.00002\Omega
	 \end{equation}
\end{Large}

\section {Incertezza su grandezze elettriche, metodo di confronto delle cadute di tensione}
È stato d'apprima impostato il multimetro come voltmetro e successivamente sono state misurate le tensioni sul provino e sul campione:
Prendiamo l'incertezza che ci riguarda per il multimetro:

\begin{Large} 
	\begin{equation}
		U_{V} = [\pm 0.0050\%rdg \pm 0.0035\% FSO]
	 \end{equation}
\end{Large}
E quindi si ottiene:
\begin{Large} 
	\begin{equation}
		V_X = 0.817 \pm 0.004 mV
	 \end{equation}
	 \begin{equation}
		V_C = 0.503 \pm 0.004 mV
	 \end{equation}
\end{Large}
\subsection{Calcolo del rapporto delle tensioni}
\begin{Large}
	\begin{equation}
		r = {V_X}/{V_C}= 1.624
	 \end{equation}
\end{Large}
\begin{Large} 
	 \begin{equation}
		U_{r} = (\frac{U_{V_X}}{V_X}+\frac{U_{V_C}}{V_C})r = 0.02086
	 \end{equation}
\end{Large}
Otteniamo:
\begin{Large} 
	\begin{equation}
		r = 1.62 \pm 0.02
	 \end{equation}
\end{Large}
\subsection{Calcolo del valore della resistenza del provino}
\begin{Large} 
	\begin{equation}
		R_{X_{CdT}} = R_c* r = 0.1624 m\Omega
	 \end{equation}
	 \begin{equation}
		U_{R_{X_{CdT}}} =(\frac{U_{R_c}}{R_c}+ \frac{U_r}{r})R_{X_{CdT}}= 0.00201 m\Omega
	 \end{equation}
	 \begin{equation}
		R_{X_{CdT}} = 0.162 \pm 0.002  m\Omega
	 \end{equation}
\end{Large}

\subsection{Calcolo della resistività del rame}
\begin{Large} 
	\begin{equation}
		\rho_{X_{CdT}} = R_{X_{CdT}}*\frac{S}{l_V} = 1,489*10^{-8} \Omega*m
	 \end{equation}
	 \begin{equation}
		U_{\rho_{X_{CdT}}} = \frac{S}{l_V}*U_{R_{X_{CdT}}}+\frac{R_{X_{CdT}}}{l_V}*U_S+\frac{S*R_{X_{CdT}}}{{l_V}^2}*U_{l_V}
	 \end{equation}
	 \begin{equation}
		\rho_{X_{CdT}} = (1.489\pm 0.001)*10^{-8} \Omega*m
	 \end{equation}
\end{Large}
\end{document}